\documentclass[11pt]{article}

\usepackage[utf8]{inputenc}
\usepackage[T1]{fontenc}
\usepackage{lmodern}
\usepackage{geometry}
\usepackage{amsmath,amssymb,amsfonts,amsthm,mathtools}
\usepackage{bm}
\usepackage{enumitem}
\usepackage{microtype}
\usepackage{hyperref}
\usepackage{titlesec}
\usepackage{booktabs}
\usepackage{array}
\usepackage{setspace}
\usepackage{mathrsfs}
\usepackage{physics}

\geometry{margin=1in}
\hypersetup{colorlinks=true, linkcolor=black, urlcolor=blue, citecolor=black}
\setlist[enumerate]{topsep=0.4em,itemsep=0.35em}
\setlist[itemize]{topsep=0.3em,itemsep=0.25em}
\titleformat{\section}{\large\bfseries}{\thesection.}{0.5em}{}
\titleformat{\subsection}{\normalsize\bfseries}{\thesubsection.}{0.5em}{}
\setstretch{1.06}

\newcommand{\Aether}{\AE{}ther}
\newcommand{\Aflow}{\AE{}ther-flow}
\newcommand{\TheoryName}{\Aflow{} Interpretation of Relativity}
\newcommand{\TheoryProgram}{\Aether{} / \Aflow{} framework}

\newcommand{\CompletedMetric}{g^{\sharp}}
\newcommand{\MatterSpace}{\Psi}

\newcommand{\Car}{\mathrm{car}}
\newcommand{\Cur}{\mathrm{cur}}
\newcommand{\Hom}{\mathrm{hom}}

\newcommand{\ForSpace}[1]{\mathcal I_{#1}^{\mathrm{for}}}
\newcommand{\PrimShell}{\mathcal S_{\mathrm{prim}}}
\newcommand{\MetricOut}{\mathscr G_{\mathrm{out}}}
\newcommand{\MatterOut}{\mathscr T_{\mathrm{out}}}

\newcommand{\ProtoSectionCore}{\mathfrak P^{\mathrm{sub,proto,sec}}}
\newcommand{\ProtoSectionPackage}{\mathfrak Q^{\mathrm{sub,proto,secgen}}}
\newcommand{\ProtoSectionState}[1]{\mathcal Q_{#1}^{\mathrm{psec}}}
\newcommand{\ProtoSectionBody}[1]{\mathcal B_{#1}^{\mathrm{psec}}}
\newcommand{\ProtoSectionFunctional}[1]{\mathscr E_{#1}^{\mathrm{psec}}}
\newcommand{\ProtoSectionShell}[1]{\mathcal Z_{#1}^{\mathrm{psec}}}

\newcommand{\PreProtoSectionPackage}{\mathfrak R^{\mathrm{sub,proto,presec}}}
\newcommand{\PreProtoSectionState}[1]{\mathcal R_{#1}^{\mathrm{ppsec}}}
\newcommand{\PreProtoSectionBody}[1]{\mathcal B_{#1}^{\mathrm{ppsec}}}
\newcommand{\PreProtoSectionProj}[1]{\Pi_{#1}^{\mathrm{ppsec}}}
\newcommand{\PreProtoSectionFiber}[2]{\mathcal F_{#1}^{\mathrm{ppsec}}(#2)}
\newcommand{\PreProtoSectionProtoMin}[1]{\mathfrak n_{#1}^{\mathrm{ppsec}}}
\newcommand{\PreProtoSectionFunctional}[1]{\mathscr F_{#1}^{\mathrm{ppsec}}}
\newcommand{\PreProtoSectionShell}[1]{\mathcal Z_{#1}^{\mathrm{ppsec}}}

\newcommand{\PreProtoSectionImageCore}{\mathfrak R^{\mathrm{sub,proto,presec,img}}}
\newcommand{\PreProtoSectionMinSectionCore}{\mathfrak R^{\mathrm{sub,proto,presec,sec}}}
\newcommand{\PreProtoSectionMinSection}[1]{\mathfrak s_{#1}^{\mathrm{ppsec,sec}}}

\newcommand{\PreProtoSectionVertical}[1]{\mathcal V_{#1}^{\mathrm{ppsec}}}
\newcommand{\PreProtoSectionResidualBody}[1]{\mathcal D_{#1}^{\mathrm{ppsec,res}}}
\newcommand{\PreProtoSectionResidualFunctional}[1]{\mathscr R_{#1}^{\mathrm{ppsec,res}}}
\newcommand{\PreProtoSectionResidualPackage}{\mathfrak R^{\mathrm{sub,proto,presec,res}}}
\newcommand{\PreProtoSectionNormalizer}[1]{\Xi_{#1}^{\mathrm{ppsec}}}

\newtheorem{definition}{Definition}
\newtheorem{proposition}{Proposition}
\newtheorem{remark}{Remark}
\newtheorem{corollary}{Corollary}
\newtheorem{theorem}{Theorem}

\title{\textbf{The \TheoryName{}}\\[0.4em]
\large Primitive-Reservoir Observer-Reduction\\
\large Proto-Section-Generating Substrate Pre-Proto-Section Dynamics\\
\large Section-Centered Normal Form Necessity / Uniqueness Theorem}
\author{}
\date{}

\begin{document}

\maketitle

\begin{abstract}
The current primitive-reservoir observer-reduction frontier already moved the live burden below the observer-relevant pre-proto-section minimizer-section core: the active theorem deriving that core from proto-section-generating substrate pre-proto-section dynamics removed the section core itself as the primitive stopping point and renewed the burden at the current pre-proto-section package. At that renewed frontier, however, a clean derivation from still more primitive explicit substrate structure is not yet in hand. The honest fallback therefore has to be sharper and internal to the current package rather than another same-output deeper relay.

The present note proves such an internal sharpening. Relative to the already fixed proto-section benchmark base and the canonical minimizer section already generated by the current pre-proto-section dynamics, every benchmark-faithful proto-section-generating substrate pre-proto-section dynamics package has a canonical section-centered residual normal form. In that normal form the package is expressed by the minimizer section, a vertical residual body over the fixed proto-section body, a nonnegative normalized residual functional vanishing exactly on the zero vertical slice, and the induced exact-shell restriction. The ambient affine presentation is no longer additional benchmark-facing freedom. It is uniquely recovered from the section-centered residual normal form by the canonical reconstruction map. Consequently the current pre-proto-section dynamics package is necessary and unique up to section-faithful normal-form equivalence once the current benchmark base and canonical section are fixed.

The benchmark boundary remains fixed throughout. The one operative metric is still exactly \(\CompletedMetric\), the ordinary matter route is unchanged, there is no second operative metric, the low-energy matter law is not enlarged, and no stronger promotion language is licensed. The positive result is therefore a genuine benchmark-facing gain strictly inside the renewed pre-proto-section package: the remaining honest burden is no longer arbitrary presentational freedom of that package, but either a derivation of its section-centered residual normal form from still more primitive explicit substrate structure or a still smaller collapse strictly inside that residual normal form.
\end{abstract}

\section{Introduction}

The active derivational program of \TheoryName{} again reaches a point where depth alone is not enough. The current active-primary theorem already showed that the observer-relevant pre-proto-section minimizer-section core is not the deepest unexplained datum on the active line, because it is canonically generated by the current proto-section-generating substrate pre-proto-section dynamics package. That theorem therefore spent the old section-core stopping point. The live burden moved to the deeper package itself.

At that renewed frontier there are only two honest options: derive or justify the current package from still more primitive explicit substrate structure, or prove a theorem-level sharpening strictly inside the current package. The present note takes the second route because a clean deeper derivation is not yet available without inventing new hidden structure. Its claim is narrower and more disciplined. It shows that, once the already generated canonical minimizer section is treated as fixed and the proto-section benchmark base is held fixed, the remaining benchmark-facing content of the current pre-proto-section package has a unique section-centered residual normal form. Arbitrary ambient affine presentation is not additional benchmark-facing freedom.

\paragraph{Framework and claim boundary.}
\TheoryProgram{} denotes the broader \Aether{} / \Aflow{} framework built on the claims that the \Aether{} is the underlying four-dimensional substrate of reality, the \Aflow{} is its intrinsic ordered motion, observed three-dimensional space is the local experiential slice of that deeper substrate, S-time is the experienced order of change arising from matter, light, and the \Aflow{}, and observed expansion is the three-dimensional appearance of deeper four-dimensional ordered motion. Within that framework, \TheoryName{} denotes the exact-closure theory in which the effective gravitational dynamics are adopted to be exactly Einsteinian with universal matter coupling. In the active sequence, \emph{adoption} means use of that established relativistic dynamics without claiming substrate derivation, while \emph{derivation} is reserved for a first-principles recovery from explicit substrate variables.

The present note stays strictly inside that boundary. It does not introduce a second operative metric, it does not enlarge the low-energy matter law, and it does not reopen the already-spent minimizer-image or larger full pre-proto-section presentations as if they were again independently live. Its narrower benchmark-facing role is to remove residual presentational freedom inside the currently live pre-proto-section package by proving a canonical section-centered normal form and a uniqueness theorem for that form.

\section{Fixed Inputs}

\begin{definition}[Recorded primitive continuation data]
\label{def:fixed}
Fix the following continuation data from the active primitive-reservoir observer-reduction line.
\begin{enumerate}
    \item For each sector label \(\sigma\in\{\Car,\Cur,\Hom\}\), a primitive forcing shell
    \[
    \ForSpace{\sigma}.
    \]
    \item The product shell
    \[
    \PrimShell
    \equiv
    \ForSpace{\Car}\times \ForSpace{\Cur}\times \ForSpace{\Hom}.
    \]
    \item The recorded metric and ordinary-matter outputs
    \[
    \MetricOut:\PrimShell\to\{\text{observer metrics}\},
    \qquad
    \MatterOut:\PrimShell\times\MatterSpace\to\{\text{observer stress tensors}\},
    \]
    satisfying
    \begin{equation}
    \MetricOut(\mathbf w)=\CompletedMetric,
    \qquad
    \MatterOut(\mathbf w,\psi)
    =
    T_{\mu\nu}^{\mathrm{matter}}[\CompletedMetric,\psi]
    \label{eq:fixedoutputs}
    \end{equation}
    for every \(\mathbf w\in\PrimShell\) and every matter configuration \(\psi\in\MatterSpace\).
\end{enumerate}
\end{definition}

\begin{definition}[Recorded proto-section benchmark base and downstream chain]
\label{def:downstream}
Fix \(\sigma\in\{\Car,\Cur,\Hom\}\). The active line already fixes the following proto-section benchmark base and downstream consequences:
\begin{enumerate}
    \item the current proto-section benchmark base \(\ProtoSectionPackage\), presented at current scope through
    \[
    \ProtoSectionState{\sigma},\qquad
    \ProtoSectionBody{\sigma},\qquad
    \ProtoSectionFunctional{\sigma},\qquad
    \ProtoSectionShell{\sigma};
    \]
    \item the current observer-relevant pre-proto-section minimizer-image core \(\PreProtoSectionImageCore\), the current observer-relevant pre-proto-section minimizer-section core \(\PreProtoSectionMinSectionCore\), and the current germ-anchored exact-shell transport-section core \(\ProtoSectionCore\);
    \item the previously recorded fact that, once the current observer-relevant pre-proto-section minimizer-section core over the fixed proto-section benchmark base is available, the current minimizer-image core, the current proto-section package, the current germ-anchored exact-shell transport-section core, and the previously recorded seed, root, foundation, ground, origin, precursor, lift, shell-restriction, split-core, selector-law, combined-response, ambient-package, trace-core, exact-shell-affine-nucleus, continuity-Hessian compressed core, and benchmark one-metric observer-package consequences on the active line follow unchanged;
    \item benchmark faithfulness, meaning preservation of the benchmark outputs \eqref{eq:fixedoutputs}, no second operative metric, no enlarged low-energy matter law, no change to the ordinary matter route, and no premature promotion from adoption to completed derivation.
\end{enumerate}
\end{definition}

\begin{definition}[Recorded proto-section-generating substrate pre-proto-section dynamics]
\label{def:preprotosec}
Fix \(\sigma\in\{\Car,\Cur,\Hom\}\). The recorded current pre-proto-section dynamics package \(\PreProtoSectionPackage\) for sector \(\sigma\) consists of the following data.
\begin{enumerate}
    \item A finite-dimensional Euclidean pre-proto-section state space \(\PreProtoSectionState{\sigma}\), a compact convex pre-proto-section body
    \[
    \PreProtoSectionBody{\sigma}\subset\PreProtoSectionState{\sigma}
    \]
    with nonempty interior, and an affine surjection
    \[
    \PreProtoSectionProj{\sigma}:\PreProtoSectionState{\sigma}\to\ProtoSectionState{\sigma}
    \]
    satisfying
    \[
    \PreProtoSectionProj{\sigma}\bigl(\PreProtoSectionBody{\sigma}\bigr)=\ProtoSectionBody{\sigma}.
    \]
    For each \(y\in\ProtoSectionBody{\sigma}\), write
    \[
    \PreProtoSectionFiber{\sigma}{y}
    \equiv
    \bigl(\PreProtoSectionProj{\sigma}\bigr)^{-1}(y)\cap\PreProtoSectionBody{\sigma}.
    \]
    Each such fiber is required to be nonempty, compact, and convex.
    \item A nonnegative smooth pre-proto-section functional
    \[
    \PreProtoSectionFunctional{\sigma}:\PreProtoSectionState{\sigma}\to[0,\infty)
    \]
    such that, for every \(y\in\ProtoSectionBody{\sigma}\), the restriction of \(\PreProtoSectionFunctional{\sigma}\) to \(\PreProtoSectionFiber{\sigma}{y}\) is strictly convex and attains a unique minimizer. The resulting minimizer depends smoothly on \(y\); write it as
    \[
    \PreProtoSectionProtoMin{\sigma}:\ProtoSectionBody{\sigma}\to\PreProtoSectionBody{\sigma}.
    \]
    These data satisfy
    \[
    \PreProtoSectionProj{\sigma}\circ\PreProtoSectionProtoMin{\sigma}
    =
    \mathrm{id}_{\ProtoSectionBody{\sigma}},
    \]
    and the effective minimum value recovers the recorded proto-section functional:
    \[
    \ProtoSectionFunctional{\sigma}(y)
    =
    \PreProtoSectionFunctional{\sigma}\bigl(\PreProtoSectionProtoMin{\sigma}(y)\bigr)
    =
    \min_{u\in\PreProtoSectionFiber{\sigma}{y}}
    \PreProtoSectionFunctional{\sigma}(u)
    \qquad
    \text{for every }
    y\in\ProtoSectionBody{\sigma}.
    \]
    \item A closed embedded exact pre-proto-section shell
    \[
    \PreProtoSectionShell{\sigma}\subset\PreProtoSectionBody{\sigma}
    \]
    satisfying
    \[
    \PreProtoSectionShell{\sigma}
    =
    \left\{
    u\in\PreProtoSectionBody{\sigma}:
    \begin{array}{l}
    \PreProtoSectionProj{\sigma}(u)\in\ProtoSectionShell{\sigma},\\[0.2em]
    \PreProtoSectionFunctional{\sigma}(u)
    =
    \displaystyle\min_{v\in\PreProtoSectionFiber{\sigma}{\PreProtoSectionProj{\sigma}(u)}}
    \PreProtoSectionFunctional{\sigma}(v)
    \end{array}
    \right\},
    \]
    and
    \[
    \PreProtoSectionProj{\sigma}\big|_{\PreProtoSectionShell{\sigma}}
    :
    \PreProtoSectionShell{\sigma}
    \xrightarrow{\cong}
    \ProtoSectionShell{\sigma}
    \]
    is a diffeomorphism.
    \item Benchmark faithfulness, meaning preservation of the benchmark outputs \eqref{eq:fixedoutputs}, no second operative metric, no enlarged low-energy matter law, and presentation as a deeper substrate-side source for the current proto-section benchmark base rather than as a replacement benchmark theory.
\end{enumerate}
\end{definition}

\begin{remark}
Definitions \ref{def:fixed}--\ref{def:preprotosec} are fixed inputs. The present note does not attempt to derive a still deeper package. It instead sharpens what is already forced inside the current pre-proto-section package once the canonical minimizer section generated there is treated honestly.
\end{remark}

\section{Section-Centered Residual Normal Form}

\begin{definition}[Canonical minimizer section and vertical model]
\label{def:vertical}
Fix \(\sigma\in\{\Car,\Cur,\Hom\}\). Let
\[
\PreProtoSectionMinSection{\sigma}
\equiv
\PreProtoSectionProtoMin{\sigma}
:
\ProtoSectionBody{\sigma}\to\PreProtoSectionBody{\sigma}
\]
be the canonical minimizer section generated by Definition \ref{def:preprotosec}(2). Let
\[
\PreProtoSectionVertical{\sigma}
\equiv
\ker D\PreProtoSectionProj{\sigma},
\]
where \(D\PreProtoSectionProj{\sigma}\) is the constant linear part of the affine map \(\PreProtoSectionProj{\sigma}\).
\end{definition}

\begin{definition}[Section-centered residual package]
\label{def:residual}
Fix a benchmark-faithful pre-proto-section package in the sense of Definition \ref{def:preprotosec}. For each \(\sigma\in\{\Car,\Cur,\Hom\}\), define the \emph{section-centered residual package} \(\PreProtoSectionResidualPackage\) by the following data:
\begin{enumerate}
    \item the fixed proto-section benchmark base \(\ProtoSectionBody{\sigma}\), \(\ProtoSectionFunctional{\sigma}\), and \(\ProtoSectionShell{\sigma}\);
    \item the canonical minimizer section \(\PreProtoSectionMinSection{\sigma}\);
    \item the residual body
    \[
    \PreProtoSectionResidualBody{\sigma}
    \equiv
    \left\{
    (y,v)\in\ProtoSectionBody{\sigma}\times\PreProtoSectionVertical{\sigma}:
    \PreProtoSectionMinSection{\sigma}(y)+v\in\PreProtoSectionBody{\sigma}
    \right\};
    \]
    \item the residual functional
    \[
    \PreProtoSectionResidualFunctional{\sigma}(y,v)
    \equiv
    \PreProtoSectionFunctional{\sigma}\bigl(\PreProtoSectionMinSection{\sigma}(y)+v\bigr)
    -
    \ProtoSectionFunctional{\sigma}(y);
    \]
    \item the residual shell
    \[
    \ProtoSectionShell{\sigma}\times\{0\}
    \subset
    \PreProtoSectionResidualBody{\sigma};
    \]
    \item benchmark faithfulness inherited from Definition \ref{def:preprotosec}(4).
\end{enumerate}
Define the reconstruction map
\[
\PreProtoSectionNormalizer{\sigma}:
\PreProtoSectionResidualBody{\sigma}\to\PreProtoSectionBody{\sigma},
\qquad
\PreProtoSectionNormalizer{\sigma}(y,v)
\equiv
\PreProtoSectionMinSection{\sigma}(y)+v.
\]
\end{definition}

\begin{proposition}[Canonical section-centered decomposition]
\label{prop:decomp}
For each \(\sigma\in\{\Car,\Cur,\Hom\}\), every \(u\in\PreProtoSectionBody{\sigma}\) can be written uniquely as
\[
u=\PreProtoSectionMinSection{\sigma}(y)+v
\]
with \(y\in\ProtoSectionBody{\sigma}\) and \(v\in\PreProtoSectionVertical{\sigma}\). Equivalently, \(\PreProtoSectionNormalizer{\sigma}\) is a bijection from \(\PreProtoSectionResidualBody{\sigma}\) onto \(\PreProtoSectionBody{\sigma}\).
\end{proposition}

\begin{proof}
Set \(y\equiv\PreProtoSectionProj{\sigma}(u)\). Then \(y\in\ProtoSectionBody{\sigma}\) by Definition \ref{def:preprotosec}(1). Put
\[
v\equiv u-\PreProtoSectionMinSection{\sigma}(y).
\]
Because \(\PreProtoSectionProj{\sigma}\) and \(\PreProtoSectionMinSection{\sigma}\) have the same value \(y\) under \(\PreProtoSectionProj{\sigma}\), the difference \(v\) lies in \(\ker D\PreProtoSectionProj{\sigma}=\PreProtoSectionVertical{\sigma}\). Hence \((y,v)\in\PreProtoSectionResidualBody{\sigma}\) and \(u=\PreProtoSectionNormalizer{\sigma}(y,v)\).

If also
\[
u=\PreProtoSectionMinSection{\sigma}(y')+v'
\]
with \(v'\in\PreProtoSectionVertical{\sigma}\), then applying \(\PreProtoSectionProj{\sigma}\) gives
\[
y'
=
\PreProtoSectionProj{\sigma}\bigl(\PreProtoSectionMinSection{\sigma}(y')+v'\bigr)
=
y.
\]
Therefore \(v'=u-\PreProtoSectionMinSection{\sigma}(y)=v\). So the decomposition is unique.
\end{proof}

\begin{proposition}[Residual normal form properties]
\label{prop:residualprops}
For each \(\sigma\in\{\Car,\Cur,\Hom\}\), the section-centered residual package satisfies:
\begin{enumerate}
    \item \(\PreProtoSectionResidualFunctional{\sigma}(y,v)\ge 0\) for every \((y,v)\in\PreProtoSectionResidualBody{\sigma}\);
    \item for each fixed \(y\in\ProtoSectionBody{\sigma}\), the unique zero of \(v\mapsto \PreProtoSectionResidualFunctional{\sigma}(y,v)\) on the fiber
    \[
    \{v\in\PreProtoSectionVertical{\sigma}:(y,v)\in\PreProtoSectionResidualBody{\sigma}\}
    \]
    is \(v=0\);
    \item the residual shell is exactly \(\ProtoSectionShell{\sigma}\times\{0\}\);
    \item under the bijection \(\PreProtoSectionNormalizer{\sigma}\), the original pre-proto-section package is recovered exactly from the residual package.
\end{enumerate}
\end{proposition}

\begin{proof}
Fix \(y\in\ProtoSectionBody{\sigma}\) and \(v\) with \((y,v)\in\PreProtoSectionResidualBody{\sigma}\). Then
\[
\PreProtoSectionMinSection{\sigma}(y)+v\in\PreProtoSectionFiber{\sigma}{y}.
\]
By Definition \ref{def:preprotosec}(2),
\[
\PreProtoSectionFunctional{\sigma}\bigl(\PreProtoSectionMinSection{\sigma}(y)+v\bigr)
\ge
\PreProtoSectionFunctional{\sigma}\bigl(\PreProtoSectionMinSection{\sigma}(y)\bigr)
=
\ProtoSectionFunctional{\sigma}(y),
\]
which proves (1).

Equality holds if and only if \(\PreProtoSectionMinSection{\sigma}(y)+v\) is also a minimizer on \(\PreProtoSectionFiber{\sigma}{y}\). Since that minimizer is unique, one must have
\[
\PreProtoSectionMinSection{\sigma}(y)+v=\PreProtoSectionMinSection{\sigma}(y),
\]
hence \(v=0\). This proves (2).

By Definition \ref{def:preprotosec}(3), an element \(u\in\PreProtoSectionBody{\sigma}\) lies in \(\PreProtoSectionShell{\sigma}\) exactly when \(\PreProtoSectionProj{\sigma}(u)\in\ProtoSectionShell{\sigma}\) and \(u\) minimizes \(\PreProtoSectionFunctional{\sigma}\) on its fiber. Under the unique decomposition of Proposition \ref{prop:decomp}, this is equivalent to \(y\in\ProtoSectionShell{\sigma}\) and \(v=0\), proving (3).

Item (4) is immediate from Proposition \ref{prop:decomp} and the definitions: the body, the shell, and the functional are exactly the pullback of the original data under \(\PreProtoSectionNormalizer{\sigma}\), and the minimizer section corresponds to the zero residual slice \(v=0\).
\end{proof}

\begin{definition}[Section-faithful normal-form equivalence]
\label{def:equiv}
Two benchmark-faithful proto-section-generating substrate pre-proto-section dynamics packages over the same fixed proto-section benchmark base are \emph{section-faithfully normal-form equivalent} if, for each sector \(\sigma\in\{\Car,\Cur,\Hom\}\), they induce the same section-centered residual package of Definition \ref{def:residual}.
\end{definition}

\begin{theorem}[Section-centered normal form necessity and uniqueness]
\label{thm:main}
Fix the primitive continuation data of Definition \ref{def:fixed}, the recorded proto-section benchmark base and downstream chain of Definition \ref{def:downstream}, and a benchmark-faithful proto-section-generating substrate pre-proto-section dynamics package in the sense of Definition \ref{def:preprotosec}. Then, for each sector \(\sigma\in\{\Car,\Cur,\Hom\}\):
\begin{enumerate}
    \item the current pre-proto-section package admits the canonical section-centered residual normal form of Definition \ref{def:residual};
    \item the original package and its section-centered residual normal form are equivalent via the canonical reconstruction bijection \(\PreProtoSectionNormalizer{\sigma}\);
    \item the exact shell is necessarily the zero residual slice over the fixed proto-section shell, and the unique minimizing section is necessarily the zero residual slice over the full fixed proto-section body;
    \item any other benchmark-faithful pre-proto-section package over the same fixed proto-section benchmark base is benchmark-facingly indistinguishable from the current one if and only if it is section-faithfully normal-form equivalent to it.
\end{enumerate}
Consequently, once the fixed proto-section benchmark base and the canonical minimizer section are held fixed, the current proto-section-generating substrate pre-proto-section dynamics package carries no residual benchmark-facing freedom beyond its canonical section-centered residual normal form. In this sense the package is necessary and unique up to section-faithful normal-form equivalence.
\end{theorem}

\begin{proof}
Items (1) and (2) are Proposition \ref{prop:decomp} together with Proposition \ref{prop:residualprops}(4). Item (3) is Proposition \ref{prop:residualprops}(2)--(3). For item (4), if two packages induce the same section-centered residual package, then by construction they are both recovered from that same residual package by the same canonical reconstruction rule, so they are benchmark-facingly indistinguishable at current scope. Conversely, if they are benchmark-facingly indistinguishable at current scope while the fixed proto-section benchmark base is held fixed, then the current package-level freedom that could still matter benchmark-facingly must agree in section-centered coordinates. That agreement is exactly the statement that the two packages induce the same residual package. Therefore section-faithful normal-form equivalence is both necessary and sufficient for benchmark-facing indistinguishability at the present frontier.
\end{proof}

\begin{corollary}[Downstream consequences are unchanged]
\label{cor:downstream}
Once the section-centered residual normal form of Theorem \ref{thm:main} is fixed, all current downstream benchmark-facing uses already recorded on the active line follow unchanged: the current observer-relevant pre-proto-section minimizer-section core, the current observer-relevant pre-proto-section minimizer-image core, the current proto-section package, the current germ-anchored exact-shell transport-section core, and the previously recorded seed, root, foundation, ground, origin, precursor, lift, shell-restriction, split-core, selector-law, combined-response, ambient-package, trace-core, exact-shell-affine-nucleus, continuity-Hessian compressed core, and benchmark one-metric observer-package consequences remain exactly as before.
\end{corollary}

\begin{proof}
Definition \ref{def:downstream}(3) records that all current downstream benchmark-facing consequences already factor through the fixed proto-section benchmark base together with the current observer-relevant pre-proto-section minimizer-section core. By Theorem \ref{thm:main}, that core and the exact shell are necessarily recovered from the section-centered residual normal form. Therefore fixing the residual normal form is sufficient to preserve all currently recorded downstream consequences unchanged.
\end{proof}

\begin{corollary}[The benchmark package remains fixed]
\label{cor:benchmark}
Every benchmark-faithful proto-section-generating substrate pre-proto-section dynamics package and every section-centered residual normal form derived from it preserve the same benchmark one-metric package \eqref{eq:fixedoutputs}. In particular, the operative metric remains exactly \(\CompletedMetric\), the ordinary matter route is unchanged, there is no second operative metric, and the low-energy matter law is not enlarged.
\end{corollary}

\begin{proof}
This is part of benchmark faithfulness in Definitions \ref{def:downstream}(4) and \ref{def:preprotosec}(4), and the residual normal form merely repackages the same benchmark-faithful data in canonical section-centered coordinates.
\end{proof}

\section{Routing Consequence}

\begin{corollary}[What the next honest move now is]
\label{cor:next}
After Theorem \ref{thm:main}, the next honest benchmark-facing manuscript must either:
\begin{enumerate}
    \item derive or justify the section-centered residual normal form of the current proto-section-generating substrate pre-proto-section dynamics package from still more primitive explicit substrate structure in a way that changes benchmark-facing status;
    \item identify a genuinely smaller theorem-level observer-relevant datum strictly inside that residual normal form and prove a sharper collapse there;
    \item do either of the previous two while preserving the same benchmark one-metric package, the same ordinary-matter route, and the same claim boundary.
\end{enumerate}
Another same-output deeper-origin relay that merely preserves the same current section-centered residual normal form and therefore merely reproduces the same current pre-proto-section package, the same current observer-relevant pre-proto-section minimizer-section core, the same current observer-relevant pre-proto-section minimizer-image core, the same current proto-section benchmark base, and the same previously recorded downstream consequences, another reopening of the larger minimizer-image or full pre-proto-section presentations, or any move that introduces a second operative metric, enlarges the ordinary matter law, or uses stronger promotion language does not satisfy the present burden.
\end{corollary}

\begin{proof}
Theorem \ref{thm:main} shows that residual benchmark-facing freedom at the current pre-proto-section frontier has already been compressed to the section-centered residual normal form. Therefore a subsequent manuscript must improve on that sharpened datum rather than merely restating it in another same-output presentation.
\end{proof}

\section{Conclusion}

The positive result of this note is a genuine benchmark-facing sharpening strictly inside the renewed pre-proto-section package. The current proto-section-generating substrate pre-proto-section dynamics package is no longer treated as carrying arbitrary benchmark-facing presentational freedom once the fixed proto-section benchmark base and the canonical minimizer section are held fixed. It admits a canonical section-centered residual normal form, and within that class it is necessary and unique up to section-faithful normal-form equivalence.

This preserves the same benchmark one-metric package and the same claim boundary. The result does not yet provide a completed first-principles derivation of GR from substrate microphysics, and it does not license stronger promotion language. Its narrower benchmark-facing gain is that the remaining honest burden is no longer unconstrained freedom of the current pre-proto-section package itself, but either a derivation of its section-centered residual normal form from still more primitive explicit substrate structure or a still smaller collapse strictly inside that residual normal form.

\end{document}
